% TOP_PROBLEM: {"id": 4800021, "title": "Multiple ergodic averages — Problem 21", "category_id": 11, "category": {"id": 11, "name": "dynamical_systems", "display_name": "Dynamical Systems", "description": "Problems about long-term behavior of deterministic systems, Hamiltonian dynamics, and periodic orbits.", "slug": "dynamical-systems", "order_index": 11, "created_at": "2026-07-31T15:26:25.671Z"}, "status": "open", "problem_number": "AMR-047-0021", "set_id": 13}
% FIRST_POSTED: 2026-10-01
% ENTRY_KIND: statement_only
% UnsolvedMath ID 4800021; source catalogue code AMR-047-0021
% !TeX program = lualatex
% Definition and sources only; this file is not an LLM solution attempt.
\documentclass[11pt,a4paper]{article}
\usepackage[margin=25mm]{geometry}
\usepackage{fontspec}
\setmainfont{lmroman10-regular.otf}[BoldFont=lmroman10-bold.otf,
 ItalicFont=lmroman10-italic.otf,BoldItalicFont=lmroman10-bolditalic.otf]
\usepackage{amsmath,amssymb,mathtools}
\usepackage{unicode-math}
\setmathfont{latinmodern-math.otf}
\AtBeginDocument{\let\setminus\smallsetminus}
\usepackage{xurl,hyperref}
\hypersetup{colorlinks=true,urlcolor=blue}
\setcounter{secnumdepth}{0}
\setlength{\parindent}{0pt}
\setlength{\parskip}{5pt}
\setlength{\emergencystretch}{3em}
\begin{document}
\section*{Multiple ergodic averages — Problem 21}\label{4800021}
{\small UnsolvedMath ID 4800021 · AMR-047-0021 · Dynamical Systems · AMR Open Problem Lists}
\subsection{Definitions and mathematical statement}
Fix $\ell\ge1$ and distinct polynomials
$p_1,\ldots,p_\ell\in\mathbb Z[t]$ with $p_i(0)=0$.
Let $T_1,\ldots,T_\ell$ be invertible measure-preserving
transformations of a common probability space
$(X,\mathcal B,\mu)$. Unlike many multiple-recurrence
questions, no commutation relations between the
transformations are imposed.

Prove that for every measurable $A\subset X$ of
positive measure there are positive integers $m,n$
such that
\[
          \mu\bigl(A\cap T_1^{m+p_1(n)}A\cap\cdots
                             \cap T_\ell^{m+p_\ell(n)}A\bigr)>0.
\]
The same $m$ and $n$ must be used for every
transformation. For an integer $k$, the notation
$T_i^kA$ denotes the image under the corresponding
invertible iterate; values of $m+p_i(n)$ may have
either sign or be zero.

Distinctness means that no two polynomials agree
identically. Since all constant terms vanish, their
pairwise differences are nonconstant. No further
linear independence condition on the polynomial
family, and no ergodicity or mixing condition on
the system, is assumed. Both positive parameters
may depend on $A$ and on the entire input data.
The target is one positive simultaneous intersection,
not a particular limiting formula or uniform lower
bound. The additional common offset $m$ is part
of the requested assertion and cannot be removed
when formulating this noncommuting version.
\subsection{Short English statement}
Does a common extra time offset restore polynomial multiple recurrence for transformations that need not commute?
\subsection{Sources}
\begin{itemize}
\item[S1] ULAM.AI and the UnsolvedMath Contributors, supplied problems.json, record 4800021 (AMR-047-0021), AMR Open Problem Lists. Catalogue identity and source transcription; CC BY 4.0.
\url{https://huggingface.co/datasets/ulamai/UnsolvedMath}.
\item[S2] Frantzikinakis - Some open problems on multiple ergodic averages (2016), Problem 21. Original source indexed by AMR-047-0021.
\url{https://arxiv.org/abs/1103.3808}.
\end{itemize}
\end{document}
